\subsection{Direct multivariate approximation of softmax}
\label{sec:softmax_approximation}

We first evaluated whether HI-SMURF could directly reproduce the coupled
probability mapping of softmax. For this experiment, we used a joint-output
conditional probability tensor whose state-dependent output vectors are
constrained to the probability simplex. The tensor-train (TT) representation
therefore predicts all output probabilities jointly. Input scores were sorted
and mean-centred before stochastic encoding, and the output permutation was
subsequently inverted. Exact softmax provided the reference probabilities for
training and evaluation; the compiled operator used neither an exponential
approximation nor a runtime output softmax.

We trained three configurations on 16,384 independent Gaussian score vectors
of dimension 16, using 2,048 validation vectors for checkpoint selection and
4,096 separate test vectors for evaluation. Each configuration was trained
once. The validation-selected model, with $N=16$ states per finite-state
machine and TT rank $r=8$, achieved a test root-mean-square error (RMSE) of
$1.296\times10^{-2}$ and a mean row total variation distance of
$5.640\times10^{-2}$ (Fig.~\ref{fig:softmax_direct}a).
The median and 95th-percentile row total variation distances were 0.0512 and
0.1026, respectively (Fig.~\ref{fig:softmax_direct}b). Residual errors
were not uniform across the probability range: larger reference probabilities
tended to be underestimated. The state-output constraints preserved
non-negativity and normalization, with a maximum test row-sum deviation of
$1.13\times10^{-6}$ in floating-point evaluation.

Validation error decreased over the recorded optimization steps for all three
configurations (Fig.~\ref{fig:softmax_direct}c). At fixed rank $r=4$,
increasing the state count from $N=8$ to $N=16$ reduced test RMSE from 0.01426
to 0.01394, while increasing the number of trainable parameters from 2,336 to
4,672. Increasing the rank from 4 to 8 at $N=16$ further reduced RMSE to
0.01296, using 16,512 parameters (Fig.~\ref{fig:softmax_direct}d).
Thus, the latter comparison yielded a 7.0\% reduction in error for a
3.53-fold increase in parameter count. Within the configurations examined,
additional TT capacity produced a modest improvement in approximation accuracy.

We next evaluated finite-stream operation on 128 additional score vectors,
using independent Bernoulli input streams and three random sequence seeds.
The finite-state machines were initialized centrally and advanced for 256
burn-in cycles before collecting $L$ categorical output events. The output
RMSE decreased from
$(4.449\pm0.214)\times10^{-2}$ at $L=64$ to
$(1.804\pm0.085)\times10^{-2}$ at $L=1024$
(mean $\pm$ sample standard deviation across sequence seeds;
Fig.~\ref{fig:softmax_direct}e). The analytic steady-state RMSE on these
same 128 vectors was $1.397\times10^{-2}$. Relative to this analytic
output, the sampling RMSE decreased from 0.0425 to 0.0114
(Fig.~\ref{fig:softmax_direct}f). Averaging the conditional probability
vectors along the same state trajectories, without categorical emission,
gave lower sampling errors. This diagnostic separates the additional
variability introduced by output emission from the error already present in
the finite-state trajectories. The finite-stream evaluation retained
floating-point TT contractions and did not include gate-level arithmetic or
hardware timing measurements.
